\documentclass[bw]{themain}
\begin{document}

    \level{5e}

    % \sequence{Expressions\ num\'eriques}
    % \cours
    % \activites{
    %     expressions_numeriques,
    %     priorites_operatoires,
    %     astuces_operatoires,
    %     expressions_numeriques_2,
    %     ecrire_expression_numerique_longueur,
    %     ecrire_expression_numerique_aire%
    % }
    % \begin{exercices}
    %     \simplesheet{0360-c,0360-d}
    %     \simplesheet{0334-c,0334-d,0371-c}
    %     \simplesheet{0334-e}
    %     \doublesheet{0374-c,0375-c}{0372-c}
    %     \doublesheet{0380-c,0381-c}{0377-c}
    %     \doublesheet{0376-c}{0379-c}
    % \end{exercices}
    % \questionsFlash{
    %     operation_unique_1,
    %     vocabulaire_operations_1,
    %     operations_multiples_1,
    %     operations_multiples_2,
    %     operations_multiples_3,
    %     calculs_priorites_5%
    % }
    % \DSAB{%
    %     duree=50 min,
    %     consignes={Ne pas répondre sur le sujet.},
    %     calculatrice=Calculatrice autorisée%
    % }{0334/3,0383/4,0384/4,0382/6,0385/3}
    % \evaluation{
    %     format=DS,
    %     sujet=C,
    %     duree=50 min,
    %     consignes={Ne pas répondre sur le sujet.},
    %     calculatrice=Calculatrice autorisée%
    % }{0334-a/6,0383-a/6,0384-a/6,0385-a/3}

    % \sequence{Proportionnalit\'e}
    % \cours
    % \activites{%
    %     a_times_x_equal_b,
    %     linearite_multiplication,
    %     coefficient_de_proportionnalite%
    % }
    % \begin{exercices}
    %     \doublesheet{0387-c,0389-c}{0388-c}
    %     \doublesheet{0390-c,0391-c}{0392-c,0393-c}
    %     \doublesheet{0394-c}{0395-c}
    %     \doublesheet{0399-c,0399-d}{0399-e,0400-c}
    %     \doublesheet{0401-c}{0402-c}
    %     \doublesheet{0407-c,0408-c,0409-c}{0410-c}
    %     \doublesheet{0411-c,0412-c,0413-c}{0414-c,0415-c}
    %     \doublesheet{0416-c,0417-c,0418-c}{0419-c,0420-c}
    % \end{exercices}
    % \questionsFlash{
    %     traduction_phrase_expression_numerique%
    % }
    % \evaluation{%
    %     format=DM,
    %     duree=Pour le 09/10/23,
    %     consignes={Répondre directement sur la feuille.}%
    %     }{0396-c/3,0397-c/3,0398-c/4}
    % \DSAB{
    %     duree=50 min,
    %     consignes={Écrire tous les calculs effectués sur votre feuille.
    %     Faire des phrases réponses. Remplir les tableaux directement sur le sujet.},
    %     calculatrice=Calculatrice autorisée,
    %     skip page id=0424%
    % }{0422/4,0421/3,0426/4,0424/4,0417/2,0425/3}
    % \evaluation{%
    %     format=DS,
    %     sujet=C,
    %     duree=50 min,
    %     font=sf,
    %     consignes={Écrire tous les calculs effectués sur votre feuille.
    %     Faire des phrases réponses. Remplir les tableaux directement sur le sujet.},
    %     calculatrice=Calculatrice autorisée,
    %     skip page id=0417%
    % }{0422-a/4,0421-a/5,0424-a/4,0417-a/4,0425-a/3}

    % \sequence{D\'ecouverte\ des\ nombres\ relatifs}
    % \cours
    % \activites{%
    %     decouverte_relatifs,
    %     comparaison_relatifs,
    %     distance_a_zero,
    %     reperage_dans_le_plan%
    % }
    % \begin{exercices}
    %     \simplesheet{0436-c}
    %     \doublesheet{0431-c,0434-c}{0435-c}
    %     \doublesheet{0432-c,0433-c}{0430-c}
    %     \doublesheet{0437-c,0439-c}{0438-c}
    %     \doublesheet{0440-c}{0441-c}
    % \end{exercices}
    % \questionsFlash{%
    %     pourcentage_1,
    %     pourcentages_2,
    %     comparaison_relatifs_1%
    % }
    % \DSAB{
    %     duree=50 min,
    %     consignes={Répondre sur votre feuille. Répondre par des phrases.},
    %     calculatrice=interdite,
    %     skip page id=0445%
    % }{0442/3,0443/2,0444/4,0445/2,0446/4,0447/5}
    % \evaluation{%
    %     format=DS,
    %     sujet=C,
    %     duree=50 min,
    %     font=sf,
    %     consignes={Répondre sur votre feuille. Répondre par des phrases.},
    %     calculatrice=interdite,
    %     skip page id=0446-a%
    % }{0442-a/3,0443-a/2,0444-a/6,0446-a/4,0447-a/5}

    % \sequence{Triangles\ partie\ 1}
    % \cours
    % \activites{
    %     inegalite_triangulaire_plus_court_chemin,
    %     inegalite_triangulaire_application,
    %     inconstructibilite_d_un_triangle,
    %     inegalite_triangulaire_verification%
    % }
    % \begin{exercices}
    %     \doublesheet{0448-c,0449-c}{0450-c, 0451-c}
    %     \simplesheet{0452-c}
    %     \doublesheet{0453-c, 0453-d}{0453-e}
    %     \simplesheet{0454-c}
    % \end{exercices}
    % \DSAB{
    %     duree=50 min,
    %     consignes={Consigness},
    %     calculatrice=interdite,
    %     skip page id=0452%
    % }{0481/3,0482/5,0452/5,0483/4,0450/3}

    % \sequence{Calcul\ litt\'eral}
    % \cours
    % \activites{
    %     introduction_calcul_litteral,
    %     conventions_ecritures,
    %     evaluer_une_expression_litterale,
    %     tester_une_egalite,
    %     reduire_une_expression_litterale_1,
    %     reduire_une_expression_litterale_2%
    % }
    % \begin{exercices}
    %     \doublesheet{0466-c}{0467-c}
    %     \doublesheet{0468-c}{0469-c}
    %     \doublesheet{0470-c}{0471-c}
    %     \doublesheet{0472-c}{0473-c}
    %     \doublesheet{0474-c,0475-c}{0476-c}
    %     \doublesheet{0477-c,0478-c}{0479-c,0480-c}
    %     \doublesheet{0484-d,0489-d}{0490-d}
    % \end{exercices}

    % \DSAB{
    %     duree=50 min,
    %     consignes={Ne pas répondre directement sur le sujet.},
    %     calculatrice=autorisée,
    %     skip page id=0469%
    % }{0484/6,0467/4,0479/3,0469/4,0470/2,0471/1}
    % \evaluation{
    %     format=DS,
    %     sujet=C,
    %     font=sf,
    %     duree=50 min,
    %     consignes={Ne pas répondre directement sur le sujet.},
    %     calculatrice=autorisée,
    %     skip page id=0469%
    % }{0484-a/6,0467-a/4,0479-a/3,0469-a/4,0471-a/3}
    % \evaluation{
    %     format=DM,
    %     duree=À rendre le ,
    %     consignes={Toutes les résponse doivent être rédigées sur une feuille
    %     à carreau format A4. Détailler les étapes de calcul.}%
    % }{0484-c/4,0489-c/4,0490-c/2}

    % \sequence{Triangles\ partie\ 2}
    % \cours
    % \activites{
    %     demontrer1,
    %     demontrer2%
    % }
    % \begin{exercices}
    %     \doublesheet{0493-c}{0494-c}
    % \end{exercices}
    % \DSAB{
    %     duree=50 min,
    %     consignes={Ne pas répondre directement sur le sujet.},
    %     calculatrice=autorisée,
    %     skip page id=0501%
    % }{0500/2,0501/6,0502/6,0503/6}
    % \evaluation{
    %     format=DM,
    %     duree=À rendre le 29/02/2024,
    %     consignes={Toutes les résponse doivent être rédigées sur une feuille
    %     à carreau format A4. Détailler les étapes du raisonnement.},
    %     skip page id=0501-a%
    % }{0500-a/2,0501-a/6,0502-a/6,0503-a/6}

    % \sequence{Angles\ et\ parall\'elisme}
    % \cours
    % \activites{
    %     caracterisation_angulaire_du_parallelisme_1,
    %     caracterisation_angulaire_du_parallelisme_2%
    % }
    % \begin{exercices}
    %     \doublesheet{0495-c,0496-c}{0497-c}
    %     \doublesheet{0504-c,0505-c}{0506-c,0507-c}
    %     \doublesheet{0508-c}{0509-c,0510-c}
    %     \doublesheet{0511-c,0512-c}{0513-c}
    % \end{exercices}
    % \questionsFlash{
    %     vocabulaire_angles_cinquieme_1,
    %     reduire_expression_litterale_1,
    %     reduire_expression_litterale_2%
    % }
    % \DSAB{
    %     duree=50 min,
    %     consignes={Répondre sur votre feuille, vous pouvez annoter le sujet.},
    %     calculatrice=autorisée,
    %     skip page id=0522%
    % }{0521/5,0522/3,0495/2,0523/5,0524/5}
    % \evaluation{
    %     format=DS,
    %     sujet = C,
    %     font = sf,
    %     consignes={Répondre sur votre feuille, vous pouvez annoter le sujet.},
    %     skip page id=0522-a%
    % }{0521-a/5,0522-a/6,0523-a/4,0524-a/5}
    % \sequence{Triangles\ partie\ 3}
    % \cours
    % \activites{
    %     somme_angles_triangle%
    % }
    % \begin{exercices}
    %     \doublesheet{0527-c,0530-c}{0528-c,0529-c,0531-c}
    % \end{exercices}
    % \questionsFlash{
    %     pourcentages_3,
    %     angle_inconnu_triangle_1%
    % }
    % \sequence{Division\ euclidienne}
    % \cours
    % \activites{
    %     multiples_diviseurs_schematisation%
    % }
    % \begin{exercices}
    %     \simplesheet{0532-c,0533-c}
    %     \doublesheet{0538-c,0539-c}{0540-c,0541-c,0542-c}
    % \end{exercices}
    % \questionsFlash{
    %     multiplication_division_par_10_100_1000,
    %     pourcentages_4%
    % }
    % \DSAB{
    %     duree=50 min,
    %     consignes={Répondre sur votre feuille, vous pouvez annoter le sujet.},
    %     calculatrice=interdite%
    % }{0550/3,0551/3,0552/2,0553/6,0554/3,0555/3}
    % \evaluation{
    %     format=DS,
    %     sujet = C,
    %     font = sf,
    %     consignes={Répondre sur votre feuille, vous pouvez annoter le sujet.},
    %     calculatrice=autorisée%
    % }{0550-a/6,0551-a/6,0552-a/2,0553-a/6}

    % \sequence{Somme\ et\ diff\'erence\ de\ nombres\ relatifs}
    % \cours
    % \activites{%
    %     somme_relatifs,
    %     regles_calcul_somme_relatif,
    %     soustraction_relatifs%
    %     }
    % \begin{exercices}
    %     \simplesheet{0560-c,0561-c,0562-c}
    %     \doublesheet{0563-c}{0564-c}
    %     \simplesheet{0560-d,0562-d,0565-c}
    %     \simplesheet{0566-c,0566-d,0570-c}
    %     \simplesheet{0571-c,0572-c}
    %     \simplesheet{0574-c,0575-c,0576-c,0577-c}
    % \end{exercices}
    % \questionsFlash{
    %     le_compte_est_bon_simple_4,
    %     somme_entiers_relatifs_1,
    %     difference_entiers_relatifs_1%
    % }
    % \evaluation{
    %     format=DM,
    %     duree=À rendre le 06/05/2024,
    %     consignes={Répondre directement sur la feuille.}
    % }{0567-c/4,0568-c/3,0569-c/4}
    % \evaluation{
    %     format=DS,
    %     sujet=A,
    %     calculatrice= interdite,
    %     duree = 50 min,
    %     consignes={Ne pas répondre sur le sujet et détailler les étapes
    %     de calcul.},
    %     skip page id=0577-d%
    % }{0563-d/4,0573-d/3,0574-d/3,0575-d/3,0576-d/3,0577-d/3,0578-c/3}

    % \sequence{Symetrie\ centrale}
    % \cours
    % \questionsFlash{
    %     somme_et_difference_entiers_relatifs_1%
    % }

    % \sequence{Calcul\ fractionnaire}
    % \cours
    % \activites{%
    %     simplification_de_fractions%
    % }
    % \begin{exercices}
    %     \simplesheet{0583-c,0584-c}
    % \end{exercices}

    \level{6e}
        \sequence{Aborder\ la\ g\'eometrie} 
            \cours
            \activites{%
                notation_geometrie,
                appartenance,
                programme_de_trace,
                programme_de_trace_2}
            \begin{exercices}
                \doublesheet{0361-c,0362-c}{0363-c,0364-c}
                \simplesheet{0365-c,0366-c}
                \doublesheet{0368-c,0369-c}{0367-c,0370-c}
            \end{exercices}
            \questionsFlash{
                tables_de_multiplication_2,
                tables_de_multiplication_3,
                tables_de_multiplication_1,
                notations_geometrie_1%
            }
            \evaluation{% 
                format=DM,
                duree=À rendre le 14/09/2022,
                consignes={Répondre directement sur la feuille.}
            }{0135-c/5,0136-a/5}
            \DSAB{%
                duree=50 min,
                consignes={Répondre directement sur la feuille.}
            }{0140/6,0135/5,0141/3,0142/3,0143/3}
            \evaluation{%
                format=DS,
                sujet=C,
                duree=50 min,
                consignes={Répondre directement sur la feuille.},
                font=sf%
            }{0140-b/6,0135-b/5,0141-b/6,0143-b/3}
        \sequence{Entiers\ naturels}
            \cours
            \begin{exercices}
                \doublesheet{0151-c}{0153-c,0154-c}
                \doublesheet{0403-c,0404-c}{0405-c,0406-c}
            \end{exercices}
            \DSAB{%
                duree=50 min,
                consignes={Pour chaque problème, répondre par au moins un
                calcul et une phrase réponse.}%
            }{0147/3,0148/3,0149/4,0403/4,0151/4,0152/2}
            \evaluation{%
                format=DS,
                sujet=C,
                duree=50 min,
                font=sf,
                consignes={Pour chaque problème, répondre par au moins un
                calcul et une phrase réponse.}%
            }{0147-a/4,0148-a/4,0149-a/4,0403-a/4,0151-a/4}
        \sequence{Angles}
            \cours
            \activites{%
                angle,
                notation_angles,nature_angles,
                tracer_un_angle%
                }
            \begin{exercices}
                \doublesheet{0165-c,0166-c}{0167-c}
                \doublesheet{0168-c,0169-c}{0170-c}
                \doublesheet{0173-c}{0172-c,0171-c}
                \doublesheet{0427-c,0427-d}{0428-c,0428-d,0429-c}
                \doublesheet{0176-c,0177-c}{0178-c}
            \end{exercices}
            \questionsFlash{
                calculs_astucieux_entiers_2,
                le_compte_est_bon_simple_1,
                notations_geometrie_2,
                le_compte_est_bon_simple_3,
                nature_des_angles_1%
            }
            \DSAB{
                duree=50 min,
                consignes={Répondre par des {\bfseries phrases}. Répondre sur
                le sujet uniquement pour la question {\bfseries 1)} de l'exercice 3.},
                skip page id=0173%
            }{0167/4,0173/3,0190/5,0428/3,0189/3,0188/2}
            \evaluation{
                format=DS,
                sujet=C,
                duree=50 min,
                font=sf,
                consignes={Répondre par des {\bfseries phrases}. Répondre sur
                le sujet uniquement pour la question {\bfseries 1)} de l'exercice 3.},
                skip page id=0173-a%
            }{0167-a/6,0173-a/6,0190-a/5,0428-a/3}

        \sequence{Division\ euclidienne}
            \activites{
                division_euclidienne%
            }
            \begin{exercices}
                \doublesheet{0155-c,0156-c}{0157-c,0158-c}
            \end{exercices}

        \sequence{Proportionnalit\'e}
        \activites{
            decouverte_proportionnalite,
            proportionnalite_linearite_additive,
            pourcentages,
            echelle%
        }
        \begin{exercices}
            \doublesheet{0390-c,0391-c}{0392-c,0393-c}
            \doublesheet{0407-d}{0408-d}
            \doublesheet{0399-c,0399-d}{0399-e}
            \doublesheet{0455-c,0456-c}{0458-c,0457-c}
            \doublesheet{0459-c,0460-c}{0461-c,0462-c}
            \doublesheet{0463-c,0464-c}{0465-c}
        \end{exercices}
        \questionsFlash{
            multiplication_a_trou_1,
            multiplication_a_trou_2,
            proportionnalite_1%
        }
        \DSAB{
            duree=50 min,
            consignes={Exercice 4 à faire directement sur le sujet.},
            skip page id=0407%
        }{0155/4,0156/3,0158/3,0399/2,0407/3,0456/3,0463/2}
        \evaluation{
                format=DS,
                sujet=C,
                duree=50 min,
                font=sf,
                consignes={Exercice 3 à faire directement sur le sujet.},
                skip page id=0407-a%
            }{0155-a/4,0156-a/3,0399-a/4,0407-a/3,0456-a/3,0463-a/2}

        \sequence{Droites\ perpendiculaires\ et\ parall\`eles}
        \cours
        \activites{
            perpendicularite,
            propriete_paralleles%
        }
        \sequence{Fractions}
            \cours
            \activites{%
                fractions_et_partages,
                fractions_et_partages_2,
                fractions_droite_graduee,
                fractions_decimales,
                ecriture_decimale,
                zeros_inutiles%
            }
            \begin{exercices}
                \doublesheet{0197-c,0198-c}{0199-c}
                \doublesheet{0193-c}{0193-d}
                \doublesheet{0194-c}{0195-c,0196-c}
                \doublesheet{0222-c}{0223-c}
            \end{exercices}
            \questionsFlash{
                comparaison_fractions_1,
                nature_des_angles_1,
                vocabulaire_operations_2%
            }
            \DSAB{
                duree=50 min,
                consignes={\noindent Répondre directement sur le sujet.},
                skip page id=0486,
                calculatrice=interdite%
            }{0198/2,0222/3,0194/6,0485/1,0486/1,0487/2,0488/5}
            \evaluation{
                format=DS,
                sujet=C,
                duree=50 min,
                font=sf,
                consignes={\noindent Répondre directement sur le sujet.},
                skip page id=0485-a,
                calculatrice=autorisee%
            }{0198-a/2,0222-a/3,0194-a/6,0485-a/2,0486-a/2,0488-a/5}
        \sequence{Proportionnalit\'e}
            \cours
            \DSAB{
                duree=50 min,
                consignes={Sur le sujet, écrire uniquement le nom et le prénom.}
            }{0206/4,0207/4,0208/4,0209/4,0210/2,0190/2,0211/2 (BONUS)}
            \evaluation{
                format=DS,
                sujet=C,
                duree=50 min,
                font=sf,
                consignes={Sur le sujet, écrire uniquement le nom et le prénom.}
            }{0206-a/5,0207-a/5,0209-a/5,0210-a/5,0190-a/2}
        \sequence{Polygones}
            \cours
            \activites{
                polygones,
                quadrilateres
            }
            \begin{exercices}
                \simplesheet{0216-c,0216-d,0217-c}
                \simplesheet{0216-e,0218-c}
                \doublesheet{0491-c}{0492-c}
            \end{exercices}
            \DSAB{
                duree=50 min,
                consignes={Sur le sujet, écrire uniquement le nom et le prénom
                et les réponses de l'exercice 1.},
                skip page id=0499%
            }{0498/4,0216/6,0499/6,0219/2,0221/2}
            \evaluation{
                format=DS,
                sujet=C,
                duree=50 min,
                font=sf,
                consignes={Sur le sujet, écrire uniquement le nom et le prénom.},
                skip page id=0499-a%
            }{0498-a/6,0216-a/6,0499-a/6,0221-a/2}
        \sequence{Ordre\ et\ comparaison}
        \evaluation{
                format=DM,
                duree=À rendre le 03/01/2023,
                font=sf,
                consignes={Répondre directement sur le sujet.}
        }{0222-c/6,0223-c/6,0224-c/6,0225-c/2}
        \DSAB{
            duree=50 min,
            consignes={Les figures doivent être tracées sur le sujet.
            Les autres réponses doivent être rédigées sur la feuille.},
            skip page id=0247%
        }{0222/3,0224/3,0247/4,0248/5,0249/5,0250/2 (BONUS)}
        \evaluation{
            format=DS,
            sujet=C,
            duree=50 min,
            font=sf,
            consignes={Les figures doivent être tracées sur le sujet.
            Les autres réponses doivent être rédigées sur la feuille.},
            skip page id=0247-a%
        }{0222-a/6,0224-a/6,0247-a/4,0248-a/4,0250-a/2 (BONUS)}
    \sequence{Op\'erations\ sur\ les\ nombres\ d\'ecimaux}
        \cours
        \activites{%
            multiplication_decimaux,
            multiplication_mental_decimaux,
            multiplier_par_01_001,
            division_decimale%
            }
        \begin{exercices}
            \doublesheet{0242-c,0242-d,0243-c}{0244-c,0245-c}
            \doublesheet{0007-c,0008-d,0009-a}{0010-a,0006-d}
            \doublesheet{0253-c,0254-c,0256-c}{0255-c}
            \doublesheet{0514-c}{0515-c,0516-c}
            \doublesheet{0517-c}{0518-c}
            \doublesheet{0519-c}{0520-c}
        \end{exercices}
        \questionsFlash{%
            le_compte_est_bon_simple_2,
            multiplier_par_05_et_025_1,
            multiplier_par_05_et_025_2,
            comparaison_fractions_2,
            multiplier_decimaux_1,
            multiplier_par_05_et_025_5,
            multiplication_division_par_10_100_1000,
            multiplication_division_par_10_100_1000_2%
        }
        \DSAB{
            duree=50 min,
            consignes={Ne pas répondre sur le sujet.},
            calculatrice = interdite%
        }{0525/4,0526/4,0008/4,0253/6,0006/2}
        \evaluation{
            format=DS,
            sujet=C,
            font=sf,
            duree = 50 min,
            consignes={Ne pas répondre sur le sujet.},
            calculatrice = interdite%
        }{0525-a/4,0526-a/4,0008-a/6,0253-a/6}
        
        \sequence{P\'erim\`etres}
        \cours
        \activites{
            conversions_unites_longueurs,
            perimetre_d_un_cercle%
        }
        \begin{exercices}
            \doublesheet{0311-c,0313-c}{0314-c}
            \doublesheet{0315-c,0316-c}{0317-c}
            \doublesheet{0318-c}{0319-c,0320-c}
            \doublesheet{0099-c,0101-c}{0100-c}
            \doublesheet{0112-c}{0113-c}
            \simplesheet{0114-c}
            \doublesheet{0534-c,0535-c}{0536-c,0537-c}
        \end{exercices}
        \questionsFlash{%
            diviseurs_1,
            conversion_unites_longueurs_1,
            conversion_unites_longueurs_2,
            multiplication_division_par_10_100_1000%
        }
        \DSAB{
            duree=50 min,
            consignes={Répondre sur votre feuille. Ne pas oublier les calculs
            et la phrase réponse.},
            endspace=15pt,
            skip page id=0329,
            calculatrice=autorisee%
        }{0549/3,0548/2,0311/4,0318/3,0320/3,0100/1,0537/4,0114/2 bonus}
        \evaluation{
            format=DS,
            sujet=C,
            duree=50 min,
            font=sf,
            consignes={Répondre sur votre feuille. Ne pas oublier les calculs
            et la phrase réponse.},
            endspace=15pt,
            calculatrice=autorisee%
        }{0549-a/4,0548-a/6,0318-a/3,0320-a/3,0100-a/2,0537-a/2,0114-a/2 bonus}

        \sequence{Aires}
        \activites{
            conversions_unites_aires,
            aire_rectangle,
            aire_triangle_rectangle%
            }
        \begin{exercices}
            \doublesheet{0323-c,0324-c}{}
            \doublesheet{0543-c,0544-c}{0545-c}
            \doublesheet{0546-c}{0547-c}
            \doublesheet{0556-c,0557-c}{0558-c,0559-c}
            \doublesheet{0325-c,0326-c}{0329-c,0327-c,0328-c}
            \simplesheet{0330-c}
        \end{exercices}
        \questionsFlash{
            nature_des_angles_3,
            calcul_perimetre,
            aire_rectangle_1%
        }
        \questionsFlash{%
        aire_rectangle_triangle_1%
        }
        \DSAB{
            duree=50 min,
            consignes = {Répondre sur votre feuille.},
            calculatrice = autorisee,
            skip page id=0325%
        }{0323/4,0545/2,0556/4,0325/2,0326/4,0330/4}
        \evaluation{
            duree = 50 min,
            format=DS,
            sujet=C,
            consignes = {Répondre sur votre feuille.},
            calculatrice = autorisee,
            skip page id=0325%
            }{0323-a/4,0545-a/4,0556-a/4,0325-a/4,0330-d/4}
            
            \sequence{Priorit\'es\ de\ calculs}
            \cours
        \begin{exercices}
            \simplesheet{0334-c,0335-c}
            \simplesheet{0334-d,0335-d}
            \simplesheet{0334-e,0335-e}
        \end{exercices}
        \DSAB{
            duree=50min,
            consignes={Répondre sur votre feuille.}%
        }{0334/6,0344/6,0335/4,0341/4}
        \evaluation{
            format=DS,
            sujet=C,
            duree=50min,
            consignes={Répondre sur votre feuille.}%
        }{0334-a/12,0335-a/4,0341-a/4}
        \evaluation{
            format=DS,
            sujet=A,
            duree=50 min,
            calculatrice = interdite,
            consignes={%
            {\normalfont\begin{tblr}{colspec=|*1{Q[l,m]}*1{Q[r]}*4{|Q[c,m]}|,hlines,row{1}={bg=gray!40!white},column{2}={font=\bfseries}}
                COMPETENCE&&A&B&C&D\\
                Je connais les priorités opératoires.&calculer&&&\\
                Je suis capable d'utiliser des astuces de calculs.&calculer&&&\\
                Je suis capable de traduire une série de calcul par une expression numérique.
                &représenter&&&\\
                Je redige correctement mes calculs.&communiquer&&&\\
                Je connais ma leçon.&transversale&&&\\
                Je présente correctement ma(mes) feuille(s), mon écriture est lisible.&transversale&&&\\
            \end{tblr}}\vskip1em\noindent Répondre directement sur la feuille.
            },
            skip page id=0579-c%
            }{0580-c/2,0579-c/6,0579-d/6,0581-c/3,0582-c/3}
            \evaluation{
                format=DS,
                sujet=B,
                duree=50 min,
                calculatrice = interdite,
                consignes={%
                {\normalfont\begin{tblr}{colspec=|*1{Q[l,m]}*1{Q[r]}*4{|Q[c,m]}|,hlines,row{1}={bg=gray!40!white},column{2}={font=\bfseries}}
                    COMPETENCE&&A&B&C&D\\
                    Je connais les priorités opératoires.&calculer&&&\\
                    Je suis capable d'utiliser des astuces de calculs.&calculer&&&\\
                    Je suis capable de traduire une série de calcul par une expression numérique.
                    &représenter&&&\\
                    Je redige correctement mes calculs.&communiquer&&&\\
                    Je connais ma leçon.&transversale&&&\\
                    Je présente correctement ma(mes) feuille(s), mon écriture est lisible.&transversale&&&\\
                \end{tblr}}\vskip1em\noindent Répondre directement sur la feuille.
                },
                skip page id=0579-c%
                }{0580-c/2,0579-c/6,0579-d/6,0581-c/3,0582-c/3}
                
                \sequence{M\'ediatrice\ d'un\ segment}
                \cours
                \activites{
                    mediatrice_equidistance,
                    construction_mediatrice_compas,
                    placer_centre_cercle,
                    construction_symetrique%
                    }
                    \begin{exercices}
                        \doublesheet{0067-c}{0068-c}
                        \simplesheet{0069-c}
                        \doublesheet{0073-c}{0074-c}
                        \doublesheet{0075-c}{0076-c,0077-c}
                        \doublesheet{0029-c}{0030-c}
                        \doublesheet{0034-c}{0035-c}
                        \doublesheet{0036-c}{0037-c}
                    \end{exercices}
                    \questionsFlash{%
                    calculs_priorites_2,
                    calculs_priorites_3,
                    calculs_priorites_4,
                    calculs_priorites_5,
                    axes_de_symetrie%
                    }
                    \DSAB{
                        duree=50 min,
                        consignes={Répondre sur le sujet uniquement pour l'exercice 2.},
                        skip page id=0102%
                        }{0067/4,0069/4,0102/4,0348/4,0349/4}
                        \evaluation{
                            duree=50 min,
                            format=DS,
                            sujet=C,
                            consignes={Répondre sur le sujet uniquement pour l'exercice 2.},
                            skip page id=0102-a%
                            }{0067-a/5,0069-a/5,0102-a/5,0348-a/5}

            \sequence{Sym\'etrie\ axiale}
            \cours
            \activites
                {%
                    symetrie_axiale,
                    symetrique_point_par_rapport_droite,
                    construction_symetrique,
                    construction_symetrique_figure_complexe,
                    construction_symetrique_compas,
                    propriete_de_conservation_symetrie_axiale_2%
                }
            \begin{exercices}
                \doublesheet{0041-c}{0042-c}
                \simplesheet{0028-c}
                \doublesheet{0029-c}{0030-c}
                \doublesheet{0031-c}{0032-c}
                \simplesheet{0265-c}
                \simplesheet{0266-c}
                \doublesheet{0036-c}{0035-c,0037-c}
            \end{exercices}
            \questionsFlash
                {
                    multiplier_par_05_et_025_1,
                    multiplier_par_05_et_025_2,
                    symetrique_oeil_nu,
                    elever_au_carre_1,
                    calculs_priorites_1,
                    multiplication_division_par_10_100_1000%
                }
            \evaluation{
                format=DS,
                sujet=A,
                duree=50 min,
                consignes={},
                skip page id=0285-c%
            }{0284-c/2,0285-c/4,0286-c/4,0287-c/4,0288-c/4,0289-c/2}
            
            \sequence{Division}
            \cours
            \activites{%
                a_est_un_diviseur_de_b,
                criteres_de_divisibilte%
            }
            \begin{exercices}
                \doublesheet{0274-c,0275-c,0276-c}{0277-c,0278-c,0279-c}
                \simplesheet{0280-d}
                \doublesheet{0296-c}{0297-c}
                \doublesheet{0306-c,0307-c}{0308-c}
            \end{exercices}
            \evaluation{
                format=DS,
                sujet=A,
                duree=50 min,
                consignes={Tous les exercices doivent être traités sur la feuille.}
            }{0270-c/3,0271-c/3,0272-c/3,0273-c/3,0306-a/4,0296-a/4}
            \evaluation{
                format=DS,
                sujet=B,
                duree=50 min,
                consignes={Tous les exercices doivent être traités sur la feuille.}
            }{0270-c/3,0271-c/3,0272-c/3,0273-c/3,0306-b/4,0296-b/4}
            \evaluation{
                format=DS,
                sujet=C,
                duree=50 min,
                font=sf,
                consignes={Tous les exercices doivent être traités sur la feuille.}
            }{0270-c/3,0271-c/3,0272-c/3,0273-c/3,0306-a/4,0296-a/4}
                            
        \sequence{Prori\'et\'es\ des\ droites\ parall\`eles\ et\ }
        \DSAB{
            duree=50 min,
            consignes={Ne pas répondre sur le sujet. Les exercices peuvent-être
            traités dans l'ordre souhaité.},
            skip page id=0118%
        }{0115/4,0116/4,0117/4,0118/4,0358/3,0359/3}
        \evaluation{
            duree=50 min,
            format=DS,
            sujet=C,
            consignes={Répondre directement sur le sujet.}%
        }{0116-a/5,0117-a/5,0115-a/5,0358-a/5}

        \sequence{Solides\ et\ volumes}
        \cours
    
        \sequence{Organisation\ et\ gestion\ de\ donn\'ees}
        \cours
        \evaluation{
            duree=50 min,
            format=DS,
            sujet=A,
            consignes={Répondre directement sur le sujet.},
            skip page id=0353-c%
        }{0352-c/4,0353-c/4,0357-c/4,0354-c/4,0355-c/4,0356-c/4}
        \evaluation{
            duree=50 min,
            format=DS,
            sujet=C,
            consignes={Répondre directement sur le sujet.},
            skip page id=0353-c%
        }{0352-c/4,0353-c/4,0357-c/4,0354-c/4,0356-c/4}

    % \level{3e}
    %     \sequence{Arithm\'etique}
    %         \cours\makeatletter
    %         \begin{exercices}
    %             \simplesheet{0137-a,0138-a}
    %         \end{exercices}
    %         \evaluation{%
    %             format=DS,
    %             sujet=A,
    %             duree=50 min,
    %             consignes={\noindent La présentation, et de la qualité
    %             de la rédaction comptera pour 2 points dans la note finale.}
    %         }{0162-a/5,0159-c/5,0161-c/4,0160-c/4}
    %         \evaluation{%
    %             format=DS,
    %             sujet=B,
    %             duree=50 min,
    %             consignes={La présentation, et de la qualité
    %             de la rédaction comptera pour 2 points dans la note finale.}
    %         }{0162-b/5,0159-c/5,0161-c/4,0160-c/4}
    %         \evaluation{%
    %             format=DS,
    %             sujet=C,
    %             duree=50 min,
    %             font=sf,
    %             consignes={La présentation, et de la qualité
    %             de la rédaction comptera pour 2 points dans la note finale.}
    %         }{0162-b/6,0159-c/6,0161-c/6}
    %     \sequence{Triangles\ semblables}
    %         \cours
    %         \begin{exercices}
    %             \simplesheet{0144-c}
    %             \doublesheet{0145-c}{0146-c}
    %             \doublesheet{0181-c}{0182-c}
    %         \end{exercices}
    %         \evaluation{%
    %             format=DS,
    %             sujet=A,
    %             duree=50 min,
    %             consignes={La présentation et la qualité de la rédaction
    %             comptera pour 2 points dans la note finale.},
    %             }{0186-a/6,0182-c/6,0187-c/6}
    %         \evaluation{%
    %             format=DS,
    %             sujet=B,
    %             duree=50 min,
    %             consignes={La présentation et la qualité de la rédaction
    %             comptera pour 2 points dans la note finale.},
    %             }{0186-b/6,0182-c/6,0187-c/6}
    %         \evaluation{%
    %             format=DS,
    %             sujet=C,
    %             font=sf,
    %             duree=50 min,
    %             consignes={La présentation et la qualité de la rédaction
    %             comptera pour 2 points dans la note finale.},
    %             }{0186-a/6,0182-c/6,0187-c/6}
    %     \sequence{Calcul\ litt\'eral}
    %         \cours
    %         \activites{double_distributivite}
    %         \begin{exercices}
    %             \doublesheet{0183-c,0184-c}{0185-c}
    %             \simplesheet{0191-c,0192-c}
    %         \end{exercices}
    %         \DSAB{
    %             duree=50 min,
    %             consignes={La présentation et la qualité de la rédaction
    %             comptera pour 2 points dans la note finale.}
    %         }{0191/6,0183/4,0185/4,0201/4}
    %         \evaluation{
    %             format=DS,
    %             sujet=C,
    %             duree=50 min,
    %             consignes={La présentation et la qualité de la rédaction
    %             comptera pour 2 points dans la note finale.}
    %         }{0191-a/8,0183-a/6,0201-a/4}
    %     \sequence{Statistiques}
    %     \sequence{Equations}
    %         \begin{exercices}
    %             \simplesheet{0229-c,0230-c,0231-c}
    %         \end{exercices}
    %         \evaluation{
    %             format=DS,
    %             sujet=A,
    %             duree=50 min,
    %             consignes={La présentation et la qualité de la rédaction
    %             comptera pour 2 points dans la note finale.}
    %         }{0235-c/4,0236-c/4,0230-c/5,0213-c/5}
    %         \evaluation{
    %             format=DS,
    %             sujet=A,
    %             duree=50 min,
    %             consignes={La présentation et la qualité de la rédaction
    %             comptera pour 2 points dans la note finale.}
    %         }{0235-c/4,0236-c/4,0230-c/5,0213-c/5}
    %     \sequence{Fonctions}
    %     \cours
    %     \begin{exercices}
    %         \simplesheet{0237-c,0238-c}
    %         \simplesheet{0239-c,0240-c}
    %         \simplesheet{0246-c,0241-c}
    %         \doublesheet{0251-c}{0252-c}
    %     \end{exercices}
    %     \evaluation{
    %         format=CC,
    %         sujet=A,
    %         duree=30min,
    %         consignes={Écrire directement sur le sujet.},
    %         skip page id=0260-a%
    %     }{0257-c/4,0258-a/4,0259-a/4,0260-a/4,0261-a/4}
    %     \evaluation{
    %         format=CC,
    %         sujet=B,
    %         duree=30min,
    %         consignes={Écrire directement sur le sujet.},
    %         skip page id=0260-b%
    %     }{0257-c/4,0258-b/4,0259-b/4,0260-b/4,0261-b/4}
    %     \evaluation{
    %         format=DS,
    %         sujet=A,
    %         duree=50 min,
    %         consignes={La présentation et la qualité de la rédaction
    %         comptera pour 2 points dans la note finale.},
    %         skip page id=0269-c%
    %     }{0268-c/6,0269-c/6,0267-c/6}
    %     \evaluation{
    %         format=DM,
    %         duree={À rendre le 06/03/2023},
    %         consignes={}
    %     }{0282-c/10,0283-c/10}

    %     \sequence{Trigonom\'etrie}
    %     \cours
    %     \activites{
    %         decouverte_trigonometrie%
    %     }
    %     \begin{exercices}
    %         \doublesheet{0292-c,0293-c,0294-c}{0295-c}
    %     \end{exercices}
    %     \evaluation{
    %         format=CC,
    %         sujet=A,
    %         duree=20 min,
    %         consignes={Répondre directement sur le sujet.}%
    %     }{0298-c/2,0299-a/2,0300-c/3,0301-a/3}
    %     \evaluation{
    %         format=CC,
    %         sujet=B,
    %         duree=20 min,
    %         consignes={Répondre directement sur le sujet.}%
    %     }{0298-c/2,0299-b/2,0300-c/3,0301-b/3}
    %     \evaluation{
    %         format=DS,
    %         sujet=A,
    %         duree=50 min,
    %         consignes={Présentation et qualité de la rédaction \hskip50pt/2},
    %         skip page id=0302-c%
    %     }{0304-c/3,0305-c/3,0302-c/6,0303-c/6}

    %     \sequence{R\'eciproque\ du\ th\'eor\`eme\ de\ Thal\`es}
    %     \cours
    %     \begin{exercices}
    %         \doublesheet{0309-c}{0310-c}
    %     \end{exercices}
    %     \DSAB{
    %         duree=50 min,
    %         consignes={La présentation et la qualité de la rédaction
    %         comptera pour 2 points dans la note finale.},
    %         skip page id=0310%
    %     }{0309/3,0310/3,0331/2,0332/4,0333/6}
    %     \evaluation{
    %         format=DM,
    %         duree=À rendre le 10/05/2023,
    %         consignes={Répondre sur une feuille A4.}
    %     }{0336-c/5,0337-c/5}

    %     \sequence{Fonctions\ lin\'eaires\ et\ affines}
    %     \cours
    %     \activites{
    %         fonctions_lineaires,
    %         fonctions_affines
    %     }
    %     \begin{exercices}
    %         \simplesheet{0338-c,0339-c,0340-c}
    %         \simplesheet{0338-d,0339-d,0340-d}
    %     \end{exercices}
    %     \DSAB{
    %         duree=50 min,
    %         consignes={La présentation et la qualité de la rédaction
    %         comptera pour 2 points dans la note finale.},
    %         skip page id=0346%
    %     }{0339/4,0345/4,0346/4,0347/6}
    %     \evaluation{
    %         duree=50 min,
    %         format=DS,
    %         sujet=D,
    %         consignes={La présentation et la qualité de la rédaction
    %         comptera pour 2 points dans la note finale.}%
    %     }{0339-d/4,0345-d/4,0346-a/4,0347-a/6}
\end{document}